% ICSE Research Track submission -- 10 pages main text + 2 page references.
% The CfP mandates exactly \documentclass[10pt,conference]{IEEEtran};
% no compsoc/compsocconf options, no spacing or font-size alterations.
\documentclass[10pt,conference]{IEEEtran}
\IEEEoverridecommandlockouts

% --- Symbols & Custom Marks ---
\usepackage{pifont}
\newcommand{\cmark}{\ding{51}}
\newcommand{\xmark}{\ding{55}}
\newcommand{\pmark}{\ding{109}} % partial coverage

% --- Core Packages ---
\usepackage{cite}
\usepackage{amsmath,amssymb,amsthm,amsfonts}
\usepackage[expansion=false]{microtype}
\usepackage{booktabs, multirow}
\usepackage{xcolor}
\usepackage{graphicx}
\usepackage{textcomp}

% --- Visualization ---
\usepackage{tikz}
\usetikzlibrary{arrows.meta,positioning,shapes.geometric,fit,backgrounds,calc}

% --- Algorithms & Code ---
\usepackage{algorithm}
\usepackage[noend]{algpseudocode}
\usepackage{listings}

% --- Layout & Utility ---
\usepackage{enumitem}
\usepackage{hyperref}
\usepackage{cleveref}
% Double-anonymous: keep PDF metadata free of author information.
\hypersetup{hidelinks, pdfauthor={}, pdfsubject={}, pdfkeywords={}}

% --- Theorem environments ---
\theoremstyle{definition}
\newtheorem{definition}{Definition}
\newtheorem{proposition}{Proposition}

% --- Listing style (view-metamodel / rule snippets) ---
\lstdefinelanguage{ATL}{
  morekeywords={module,create,from,rule,to,lazy,unique,helper,def,
                context,uses,thisModule,and,or,not,in},
  morekeywords=[2]{@llm,@check}, alsoletter={@},
  sensitive=true, morecomment=[l]{--}, morestring=[b]'}
\lstdefinelanguage{KM3}{
  morekeywords={package,class,attribute,reference,container,
                enumeration,literal,datatype,oppositeOf},
  sensitive=true, morecomment=[l]{--}}
\lstset{frame=single, backgroundcolor=\color{gray!5},
        basicstyle=\ttfamily\scriptsize, breaklines=true,
        showstringspaces=false, numbers=none, upquote=true,
        keywordstyle=\bfseries, keywordstyle={[2]\bfseries\color{purple}},
        commentstyle=\itshape\color{black!60},
        captionpos=t, abovecaptionskip=2pt, belowcaptionskip=5pt, framesep=2pt,
        aboveskip=6pt, belowskip=4pt, xleftmargin=2pt, xrightmargin=2pt}
\renewcommand{\lstlistingname}{Listing}
% IEEEtran's \@makecaption lets listing frames overlap the caption; use a plain one.
\makeatletter
\def\lst@makecaption#1#2{{\centering\footnotesize #1.\ #2\par}\vspace{4pt}}
\makeatother
\crefname{lstlisting}{Listing}{Listings}
\Crefname{lstlisting}{Listing}{Listings}
\crefname{algorithm}{Algorithm}{Algorithms}
% unique hyperref anchors for algorithmic line numbers across algorithms
\makeatletter\newcommand{\theHALG@line}{\thealgorithm.\arabic{ALG@line}}\makeatother

% --- Math shorthands (AgentM2M) ---
\newcommand{\tool}{\textsc{AgentM2M}}
\newcommand{\auto}{\textsc{AutoM2M}}
\newcommand{\devteam}{\textsc{DevTeam}}
\newcommand{\MM}{\mathit{MM}}
\newcommand{\Mods}[1]{\mathcal{M}(#1)}
\newcommand{\TL}{\mathit{TL}}
\newcommand{\team}{\mathcal{T}}
\newcommand{\Team}{\Theta}
\newcommand{\str}{\mathrm{str}}
\newcommand{\sem}{\mathrm{sem}}
\newcommand{\Dist}[1]{\mathcal{D}(#1)}
\newcommand{\cls}[1]{\textsf{#1}}
\newcommand{\W}[1]{\textsf{W#1}}
\newcommand{\D}[1]{\textsf{D#1}}
\newcommand{\RC}[1]{\textbf{RC#1}}
\newcommand{\code}[1]{\texttt{#1}}
\newcommand{\step}[1]{\textcircled{\scriptsize #1}}
\newcommand{\fp}{\mathit{fp}}
\newcommand{\Obl}{\mathit{Obl}}
\newcommand{\todo}[1]{\textcolor{red}{\textbf{[TODO: #1]}}}
\DeclareRobustCommand{\tracemark}{\tikz[baseline=-0.6ex]\node[draw,diamond,fill=black!70,inner sep=1.1pt]{};}
\IfFileExists{tables/numbers.tex}{\input{tables/numbers.tex}}{}

\begin{document}

\title{\textsc{AutoM2M}: Checking the Composition of\\Automatically Assembled LLM Agent Teams}

% --- Double-anonymous review: no author block in the submission. ---
% \author{...}

\maketitle

\begin{abstract}
\todo{abstract}
\end{abstract}

\begin{IEEEkeywords}
LLM-based multi-agent systems, model transformations, architectural mismatch, failure attribution, code generation
\end{IEEEkeywords}

% ===========================================================================
\section{Introduction}\label{sec:intro}
% ===========================================================================
\todo{intro}

% ===========================================================================
\section{Background}\label{sec:background}
% ===========================================================================

\auto{} builds on \tool{}~\cite{agentm2m2026}, which treats every hand-off
between two LLM agents as a model-to-model (M2M) transformation. We introduce
its concepts on one running example and state them as definitions that
\Cref{sec:approach} builds on.

\subsection{Running example}\label{sec:example}

\devteam{} builds a small to-do service. The \emph{Analyst} turns requests into
user stories with acceptance criteria; the \emph{Architect} designs one API
operation per accepted story; the \emph{Tester} writes executable tests; the
\emph{Developer} implements each operation. The requirements contain story
\code{S1} \emph{create a task} (criterion \code{S1.1}: a title is required),
\code{S2} \emph{mark a task done} (\code{S2.1}: completing a done task is
rejected) and \code{S3}, a draft. In a team built by CaptainAgent or
AutoAgents~\cite{song2024captain,chen2024autoagents}, these four agents would be
four role prompts and a plan in prose. In \tool{}, each is a form, and each
hand-off is a rule set.

\subsection{Views, models and hybrid rules}

\begin{definition}[View metamodel, model]\label{def:mm}
A \emph{view metamodel} $\MM=(C,F)$ is a set of classes $C$ and typed
features $F$: attributes with a primitive type, and references to classes
(possibly of other views), each with a lower bound (\emph{mandatory} if
$\ge 1$) and an upper bound. A \emph{model} $M\in\Mods{\MM}$ is a set of
objects that \emph{conforms} to $\MM$.
\end{definition}

Each agent owns one view, the form it fills in (\Cref{lst:views}). The
Architect's \cls{Operation} knows nothing about stories except through a
reference, so each agent stays in its own vocabulary, as in view-based
development~\cite{klare2021vitruvius,bruneliere2019views}.

\begin{lstlisting}[language=KM3, caption={View metamodels of \devteam{} (KM3-like).}, label={lst:views}]
package Req  { class UserStory { attribute id, title, status : String;
                 reference criteria[0-*] : Criterion; }
               class Criterion { attribute id, text : String;
                 reference story[1] : UserStory; } }
package Arch { class Operation { attribute name, signature : String;
                 reference story[1] : Req!UserStory; } }
package Test { class TestSuite { attribute name, code : String;
                 reference story[1] : Req!UserStory; } }
package Code { class CodeEdit  { attribute file, body : String;
                 reference operation[1] : Arch!Operation;
                 reference story[1] : Req!UserStory; } }
\end{lstlisting}

\begin{definition}[Hybrid rule]\label{def:rule}
A \emph{hybrid rule} $r$ has a source pattern (typed variables over
source views), an optional guard, and a target pattern creating one object of
class $D$. Its bindings are partitioned into \emph{structural bindings}
$f\leftarrow e$, where $e$ is a navigation path or literal evaluated by a
deterministic engine, and \emph{stochastic bindings}
$f\leftarrow\code{@llm}(p,\fp)$ with a prompt $p$, a \emph{footprint} $\fp$
(the set of navigation paths an LLM may read), and a \emph{validator}
$\mathit{chk}$ that a sampled value must pass. A \emph{hand-off} $T$ is a set
of rules with source views $\mathit{src}(T)$ and one target view
$\mathit{tgt}(T)$.
\end{definition}

\begin{lstlisting}[language=ATL, caption={Hand-off \code{Arch2Code} of \devteam{}. \code{@llm} marks a value an LLM writes; \code{@check} its validator.}, label={lst:rule}]
rule Op2Edit {
  from op : Arch!Operation, ts : Test!TestSuite (ts.story = op.story)
  to   e  : Code!CodeEdit (
       file <- op.name, operation <- op, story <- op.story,  -- structural
       body <- @llm('Implement the operation',
                    Sequence{op.signature, op.story.criteria}), -- footprint
       @check body.passes_tests(op.story, ts.code) ) }        -- validator
\end{lstlisting}

The engine computes the structure of the hand-off: for every match of the
source pattern it creates one target object, fills structural bindings and
resolves references, as ATL's matched rules do~\cite{jouault2008atl}. Only a
value no formula can compute, here an implementation, is sampled from an
LLM, which sees only the footprint. A sample is kept only if the validator
accepts it; otherwise the engine re-samples, at most $k$ times, and then
escalates. Because \code{Op2Edit} matches \emph{every} pair of an operation
and its test suite, no operation can be forgotten by an LLM.

\begin{definition}[Trace link, stamp, obligation]\label{def:trace}
When rule $r$ fires on match $m$ and creates $t$, the engine records a
\emph{trace link} $(m,t,r)$~\cite{jouault2005traceability,vara2014metagem}.
For every accepted stochastic value $t.f$ it records a \emph{stamp}, the digest
of the footprint's value at acceptance time. A value whose stamp differs
from the digest of its current footprint is an \emph{obligation}.
\end{definition}

Stamps make change propagation exact: when the Analyst edits \code{S2.1},
only the values whose footprints read \code{S2.1} become obligations. For
this paper, the important consequence is different: the complete input of
every LLM-written value is recorded and small, which is the information that
failure attribution needs~\cite{chen2026traceelephant}.

\begin{definition}[Team model, acceptance]\label{def:team}
A \emph{team model} $\team=(A,V,\mathbf{T},\omega,\varphi)$ has agents $A$,
views $V$, hand-offs $\mathbf{T}$, write rights $\omega:A\to 2^V$, and an
\emph{acceptance predicate} $\varphi$ evaluated by the engine: every match is
covered by a trace link, no obligation or escalation is open, all stamps are
fresh, and all validators pass.
\end{definition}

The team is done when $\varphi$ holds, not when an agent says so. Adding an
agent at runtime is a higher-order transformation (HOT)~\cite{tisi2009hot}
over $\team$, kept causally connected to the running team
(models@run.time~\cite{blair2009runtime}).

\subsection{What \tool{} leaves open}

\tool{}'s team model, metamodels and rules are written by a person. Its
evaluation on DevBench~\cite{li2024devbench} named ``learning the glue'' as
its first open problem. For an automatically assembled team, the team
builder is the architect: it must write the forms, the rules and $\varphi$, and
nothing guarantees that it writes good ones. The composition gap does not
disappear; it moves from the messages to the team specification. The next
section checks that specification.

% ===========================================================================
\section{Approach}\label{sec:approach}
% ===========================================================================

\auto{} lifts \tool{}'s division of labour one level up. \tool{} lets an LLM
write \emph{values} while a deterministic engine owns the \emph{structure} of
each hand-off. \auto{} lets an LLM \emph{propose a team} while a
deterministic checker owns the decision to \emph{admit} it
(\Cref{fig:pipeline}). Algorithm~\ref{alg:admit} admits a team;
Algorithm~\ref{alg:run} runs it, attributes failures and repairs it.

\begin{figure}[t]
\centering
\begin{tikzpicture}[font=\sffamily\scriptsize, >={Stealth[length=3pt]},
  llm/.style={draw=purple!70!black, fill=purple!6, rounded corners=2pt, align=center, minimum height=7mm, inner sep=2pt},
  eng/.style={draw=green!40!black, fill=green!6, rounded corners=2pt, align=center, minimum height=7mm, inner sep=2pt},
  lbl/.style={font=\sffamily\tiny, align=center}]
  \node[draw, rounded corners=2pt, minimum height=7mm, align=center] (task) at (0,0) {Task};
  \node[llm] (build) at (1.45,0) {A: Builder\\proposes $\Team$};
  \node[eng] (check) at (3.35,0) {B: Checker\\\W1--\W6};
  \node[eng] (run) at (5.2,0) {D: \tool{}\\runtime};
  \node[eng] (done) at (6.95,0) {$\varphi$ holds};
  \node[eng] (attr) at (5.2,-1.35) {E: Attribution\\trace + replay};
  \node[llm] (delta) at (2.4,-1.35) {F: Builder\\proposes $\Delta$};
  \draw[->] (task) -- (build); \draw[->] (build) -- (check);
  \draw[->] (check) -- node[above, lbl] {admit} (run);
  \draw[->] (run) -- (done);
  \draw[->] (check.north) to[bend right=30] node[above, lbl] {C: diagnostics} (build.north);
  \draw[->] (run) -- node[right, lbl] {$\varphi$ fails} (attr);
  \draw[->] (attr) -- node[above, lbl] {fault report} (delta);
  \draw[->] (delta) -| node[pos=0.3, below, lbl] {check $\Team\oplus\Delta$} (check);
\end{tikzpicture}
\caption{\auto{}. LLM steps (purple) propose; deterministic steps (green)
decide. A team runs only after the checker admits it; every repair passes the
same checker.}
\label{fig:pipeline}
\end{figure}

\subsection{Step A: typed teams}\label{sec:typed}

\begin{definition}[Typed team]\label{def:typed}
A \emph{typed team} is $\Team=(A,V,\mathbf{T},\omega,\kappa,\tau,G,\varphi)$: a
team model (\Cref{def:team}) whose views include a \emph{goal view} $\MM_0$
into which the task is lifted, extended with $\kappa:A\to 2^{\mathit{Tools}}$,
the tools of each agent; $\tau$, the tools each stochastic binding's validator
needs; and a set $G$ of \emph{goal obligations} $(C,s,k)$, a goal-view class
$C$, a scope $s$ and a kind $k\in\{\mathit{checked},\mathit{delivered}\}$.
Every validator comes from a library that declares its \emph{strength}
(\emph{form}: parses, compiles; \emph{behaviour}: executes the value against
examples or tests) and its tool needs.
\end{definition}

A builder fills a constrained JSON format, not free ATL: footprints are
navigation paths and validators are library entries. Agents keep free-text
role prompts; how an agent \emph{thinks} remains prose, but what it receives
and produces is typed. \Cref{lst:proposal} renders a first proposal for
\devteam{} that we wrote to contain one instance of each defect class.

\begin{lstlisting}[caption={Seeded proposal for \devteam{} (rendering of its JSON; defects marked).}, label={lst:proposal}]
agents: Analyst{} Architect{} Developer{exec} Tester{}        -- D5
writes: Analyst:Req Architect:Arch Developer:Code,Test Tester:Test -- D3
goal:   (Criterion, c.story.status='accepted', checked)
rule Story2Operation { from s:Req!UserStory (s.status='accepted')
   to op:Arch!Operation(name<-s.title, story<-s,
      signature<-@llm('Derive an API signature', s.title) @check compiles) }
rule Op2Edit { from op:Arch!Operation  to e:Code!CodeEdit(
      body<-@llm('Implement', {op.signature, op.returnType})     -- D2
      @check compiles) }
rule Edit2TestRun { from e:Code!CodeEdit  to r:Test!TestRun(   -- D1
      verdict<-@llm('Run and judge', e.body) @check runs) }
done:  cover(G), valid, fresh, noObl, Tester.says('ALL TESTS PASS') -- D4
\end{lstlisting}

\subsection{Step B: six well-formedness conditions}\label{sec:check}

The checker decides six conditions on $\Team$ alone, with no task data, LLM
call or execution. Two graphs support them. The \emph{type graph}
$\mathcal{G}_\Team$ has an edge $C\to D$ for every rule with $C$ in its source
pattern that creates $D$. The \emph{feature-level data-flow graph}
$\mathcal{F}_\Team$ has an edge $(C,f)\to(D,g)$ whenever a binding that writes
feature $g$ of $D$ reads feature $f$ of $C$: in its expression, its footprint,
or an argument that its validator reads (\code{op.story.criteria} reads
\code{Operation.story}, \code{UserStory.criteria} and the attributes of
\cls{Criterion}). Let $B$ be the set of features written by stochastic bindings
with a behavioural validator.

\paragraph*{\W1, hand-offs are well typed} Every rule of $T$ reads only
classes of $\mathit{src}(T)$ and creates a class of $\mathit{tgt}(T)$; every
bound feature exists and its expression type-conforms; every footprint and
validator-argument path type-checks against $\mathit{src}(T)$; stochastic
bindings target attributes and have at least one validator. \emph{Rules out
the formal part of \D2}: \code{op.returnType} in \Cref{lst:proposal} names
information no one produces.

\paragraph*{\W2, one writer per artefact} The sets $\omega(a)$ are pairwise
disjoint and cover $V\setminus\{\MM_0\}$, and each class is created by exactly
one rule. \emph{Rules out \D3}: Developer and Tester both write \cls{Test}.

\paragraph*{\W3, targets are complete} Every mandatory feature of a created
class is bound exactly once, and every structural reference binding is
\emph{resolvable}: the source class it yields equals the reference type, or
some rule maps it there. \emph{Rules out the content part of \D2}: no consumer
receives an empty mandatory slot or a dangling link.

\paragraph*{\W4, every obligation reaches a check that reads it} (a)~For each
$(C,s,\mathit{checked})\in G$, some feature of $C$ reaches $B$ in
$\mathcal{F}_\Team$ (\emph{anchored coverage}); for each
$(C,s,\mathit{delivered})\in G$, the declared deliverable feature is in $B$,
references $C$, and is reachable from a feature of $C$. (b)~Every view is
reachable from $\MM_0$ in the hand-off graph and reaches a view containing a
feature of $B$. \emph{Rules out \D1}: in \Cref{lst:proposal} no binding reads
any feature of \cls{Criterion}. A path in $\mathcal{G}_\Team$ is not enough:
a builder can create one test per criterion whose oracle never sees the
criterion (\Cref{sec:rq2static}). Guards are not decidable in general; the
checker warns when a guard on an obligation's flow differs from its scope,
and the runtime clause $\mathit{cover}(G)$ enforces coverage per instance.

\paragraph*{\W5, the engine decides completion} $\varphi$ is exactly the
engine clauses $\mathit{cover}(G)$ (every in-scope goal object has a
validated descendant, and a validated deliverable if delivered),
$\mathit{valid}$, $\mathit{fresh}$ and $\mathit{noObl}$, and the hand-off
graph is acyclic~\cite{klare2023termination}. \emph{Rules out \D4}:
\code{Tester.says(...)} is not engine-checkable.

\paragraph*{\W6, work goes to agents that can do it} For every stochastic
binding $b$ of a hand-off whose target view agent $a$ owns,
$\tau(b)\subseteq\kappa(a)$. \emph{Rules out \D5}: \code{runs} needs code
execution, which the Tester lacks.

\begin{proposition}[Cost]\label{prop:cost}
\W1--\W6 are decidable in time polynomial in $|\Team|$ (classes, features,
rules, bindings, path segments).
\end{proposition}
\noindent\emph{Sketch.} \W1 navigates each path once; \W2, \W3, \W6 compare
finite sets; \W4 is reachability in $\mathcal{F}_\Team$, whose size is bounded
by bindings times path segments, and in the hand-off graph; the cycle test of
\W5 is a topological sort.

\begin{proposition}[Admission guarantees]\label{prop:guarantee}
If $\Team$ is admitted and the runtime stops with $\varphi$ true, then (i)~every
in-scope goal object has a descendant whose value passed a behavioural
validator, and for \emph{delivered} obligations a validated deliverable,
whose inputs read data derived from that object; (ii)~no object has two
producers and no view two writers; (iii)~every LLM-written value was produced
from a recorded footprint under a stamp and passed its validator;
(iv)~completion was decided by engine clauses only.
\end{proposition}
\noindent\emph{Sketch.} (i) combines \W4(a) at type level with ATL's
matched-rule semantics and $\mathit{cover}(G)$ at instance level; (ii) is
\W2; (iii) is \W1 with the runtime; (iv) is \W5. \emph{Not guaranteed}: that
the forms capture what the task needs, that validators are strong enough, or
that agents reason correctly.

\subsection{Step C: diagnostics and admission}\label{sec:admit}

A failed condition yields a \emph{diagnostic}: the condition, the offending
element, and how to fix it. \Cref{lst:diag} is the checker's actual output on
\Cref{lst:proposal}: one violation per seeded defect and nothing else.
Diagnostics are complete for the conditions checked and name a rule, path or
write right rather than ``improve coordination''. Algorithm~\ref{alg:admit}
feeds them back to the builder for at most $k_{\mathit{adm}}$ rounds.

\begin{lstlisting}[caption={Checker output on the seeded proposal.}, label={lst:diag}]
W1 Op2Edit.body: footprint path 'op.returnType': Arch!Operation has no
   feature 'returnType' (its features are ['name','signature','story'])
W2 view Test written by ['Developer', 'Tester']; needs exactly one writer
W4 goal obligation (Criterion, ...): no behavioural validator reads data
   derived from Criterion
W5 done-clause Tester.says('ALL TESTS PASS') is not engine-checkable
W6 Edit2TestRun.verdict needs ['exec']; owner Tester has []
REJECTED: 5 violation(s)
\end{lstlisting}

\begin{algorithm}[t]
\caption{Admission (Steps A--C)}\label{alg:admit}
\small
\begin{algorithmic}[1]
\Require task $x$, builder LLM $\mathcal{B}$, budget $k_{\mathit{adm}}$
\State $\MM_0, M_0 \gets \Call{Lift}{x}$ \Comment{goal view and model, deterministic}
\State $\Team \gets \mathcal{B}.\Call{Propose}{x, \MM_0}$; $i\gets 0$
\Loop
  \State $\mathit{diag}\gets\emptyset$
  \For{$c \in \langle\W1,\dots,\W6\rangle$}
     \State $\mathit{diag}\gets\mathit{diag}\cup\Call{Violations}{c,\Team,\mathcal{G}_\Team,\mathcal{F}_\Team}$
  \EndFor
  \If{$\mathit{diag}=\emptyset$} \Return \Call{Compile}{$\Team$, $M_0$} \Comment{admitted}
  \EndIf
  \If{$i = k_{\mathit{adm}}$} \Return $\textsc{NotAdmitted}(\mathit{diag})$ \EndIf
  \State $\Team \gets \mathcal{B}.\Call{Revise}{\Team,\mathit{diag}}$; $i\gets i+1$
\EndLoop
\end{algorithmic}
\end{algorithm}

\subsection{Steps D--F: run, attribute, repair}\label{sec:runloop}

An admitted team is compiled onto the unchanged \tool{} runtime: views become
metamodels, hand-offs become rule modules with \code{@check} expressions over
the validator library, and the task is lifted into the goal model.
Algorithm~\ref{alg:run} runs it to a fixpoint and evaluates $\varphi$, then
attributes and repairs failures.

\paragraph*{Locate} A failed $\mathit{valid}$ or $\mathit{noObl}$ clause names
a pair (target object, stochastic binding); through the trace this gives the
rule, the hand-off, the owning agent and the stamped footprint. A failed
$\mathit{cover}(G)$ clause names a goal object with no validated descendant.
Failures that are consequences of a located one (the coverage gap of the same
goal object, stale stamps behind an escalation) are not attributed again.
Location is a lookup, not an inference over a transcript.

\paragraph*{Classify} Because the footprint is recorded and small, the engine
replays the binding (lines~\ref{l:val}--\ref{l:spec}): a validator that gives
two verdicts on the same value is a \emph{validator fault}; if a replay on
the same footprint passes, a \emph{sampling fault}; if a replay on the
footprint widened by one hop (every attribute of each source variable and of
the objects its references reach) passes, a \emph{footprint fault}, an
instance of \D2; if re-sampling an LLM-written value in the footprint lets the
binding pass, an \emph{upstream fault}, attributed to that upstream binding
and its agent; otherwise a \emph{specification fault}. With $r$ replays per
step the procedure makes at most $2r+2r$ LLM calls per failure.

\paragraph*{Repair} Each fault class has a remedy, and every team change is
checked before it is applied (line~\ref{l:checkdelta}): sampling faults get a
fresh budget; upstream faults turn the upstream value into an obligation;
footprint faults widen the footprint by the paths that made the replay pass;
specification and coverage faults go back to the builder, which proposes a
team delta $\Delta$. An admitted delta is applied to the \emph{running} team:
rule-level edits regenerate one module and invalidate exactly the stamps whose
prompt, validator or footprint changed; additive edits (new views, agents,
hand-offs) are a HOT; anything else rebuilds the team. Accepted values outside
the change are kept, which a free-form builder, rebuilding from scratch,
cannot do.

\begin{algorithm}[t]
\caption{Run, attribute and repair (Steps D--F)}\label{alg:run}
\small
\begin{algorithmic}[1]
\Require admitted $\Team$, goal model $M_0$, budgets $k$ (samples), $k_{\mathit{rep}}$, $r$
\For{$j = 0,\dots,k_{\mathit{rep}}$}
  \State $\mathit{fails}\gets\Call{RunToFixpoint}{\Team, M_0, k}$; \textbf{if} $\varphi$ \textbf{then return} \textsc{Done}
  \If{$j = k_{\mathit{rep}}$} \Return \textsc{Failed}(\textit{fails}) \EndIf
  \For{$f\in\Call{Roots}{\mathit{fails}}$} \Comment{drop consequential failures}
    \State $(t,b,\rho,a,\fp)\gets\Call{Locate}{f,\TL}$ \Comment{trace lookup}
    \If{$\mathit{chk}_b(v)\neq\mathit{chk}_b(v)$} $c\gets$ validator \label{l:val}
    \ElsIf{$\exists$ pass among $r$ samples on $\fp$} $c\gets$ sampling
    \ElsIf{$\exists$ pass among $r$ samples on $\Call{Widen}{\fp}$} $c\gets$ footprint
    \ElsIf{$\exists u\in\fp$ written by $b'$: re-sampling $b'$ lets $b$ pass} $c\gets$ upstream$(b')$
    \Else\ $c\gets$ specification \label{l:spec}
    \EndIf
    \State $\mathit{report}\gets\mathit{report}\cup\{(c,t,b,\rho,a)\}$
  \EndFor
  \State $\Delta\gets\Call{Remedy}{\mathit{report}}\cup\mathcal{B}.\Call{ProposeDelta}{\Team,\mathit{report}}$
  \If{$\Call{Check}{\Team\oplus\Delta}=\emptyset$} $\Team\gets\Call{ApplyInPlace}{\Team,\Delta}$ \label{l:checkdelta}
  \EndIf
\EndFor
\end{algorithmic}
\end{algorithm}

\paragraph*{What \auto{} does not solve} The checker verifies that a team fits
together, not that its forms capture the task: an essential concept missing
from the goal view is invisible to it. \W4 demands a behavioural validator but
not a strong one; weak validators turn wrong values into accepted ones. And
errors inside one agent's reasoning are outside any composition check; at
best they are classified as sampling or specification faults.

% ===========================================================================
\section{Related Work}\label{sec:related}
% ===========================================================================

\paragraph*{Automatic team construction} Team builders generate roles, plans,
state machines or workflow code per task:
CaptainAgent, AutoAgents, AgentVerse, DyLAN, EvoAgent, MetaAgent and
MAS-Zero~\cite{song2024captain,chen2024autoagents,chen2024agentverse,liu2024dylan,yuan2025evoagent,zhang2025metaagent,ke2025maszero},
and offline searches over workflows and topologies such as ADAS, AFlow,
GPTSwarm and MaAS~\cite{hu2025adas,zhang2025aflow,zhuge2024gptswarm,zhang2025maas}.
Their quality signal is an end-to-end score, an LLM critic or observer, or
execution feedback; MetaAgent's finite state machine makes control explicit
but not the content passed between states~\cite{zhang2025metaagent}, and
self-improving builders recover feedback from traces as textual
gradients~\cite{he2026tpgo}. \auto{} does not propose a new search
strategy; it proposes a target representation that any of them could emit
and a check that any of them could run inside its loop. Hand-designed
software teams (MetaGPT, ChatDev, self-collaboration)~\cite{hong2024metagpt,qian2024chatdev,dong2024selfcollab}
fix roles and hand-offs by hand; He, Treude and Lo survey the
area~\cite{he2025lma}.

\paragraph*{Failure analysis and attribution} MAST locates most multi-agent
failures in system design and inter-agent misalignment~\cite{cemri2025mast};
AgentAsk places them on hand-off edges and intervenes with clarification
questions~\cite{lin2026agentask}; Kim et al.\ show that whether a team helps
depends on whether its structure fits the task~\cite{kim2025scaling}; Lu et
al.\ report that frameworks complete about half of programmable tasks and
that communication is rarely analysed~\cite{lu2025exploring}. Who\&When and
TraceElephant attribute failures from transcripts and find agent- and
step-level attribution
unreliable~\cite{zhang2025whowhen,chen2026traceelephant}. \auto{} attacks the
same hand-offs statically and makes attribution a lookup plus a bounded
replay over recorded footprints.

\paragraph*{Structured communication} MetaGPT's documents, DSPy signatures,
schema-validated shared state (PatchBoard), Agora's protocol documents, and
MCP/A2A make individual messages well formed or reliably
transported~\cite{hong2024metagpt,khattab2024dspy,zhang2026patchboard,marro2024agora,mcp2024,a2a2025};
La Malfa et al.\ argue that LLM teams ignore classical coordination
protocols~\cite{lamalfa2025miss}, which KQML-style languages
provided~\cite{finin1994kqml}. These approaches check \emph{messages}.
\auto{} checks the \emph{team}: coverage of obligations, single writers,
complete targets, engine-decided completion and capabilities, across all
hand-offs.

\paragraph*{Software architecture and agent-oriented SE} Architectural
mismatch~\cite{garlan1995mismatch,garlan2009mismatch}, design by
contract~\cite{meyer1992dbc} and formal connectors~\cite{allen1997connection}
motivate our conditions, and Gaia's role schemas already stated
responsibilities, permissions and protocols~\cite{wooldridge2000gaia}. Ad hoc
teamwork studies agents that meet without prior
coordination~\cite{stone2010adhoc,wang2024oaht}; a team builder is in the
opposite position, since it creates all teammates and can coordinate them in
advance.

\paragraph*{Model-driven engineering} \auto{} reuses ATL's execution model
and typing~\cite{jouault2008atl}, persisted
traceability~\cite{jouault2005traceability,vara2014metagem}, incremental
execution~\cite{jouault2010incremental}, HOTs~\cite{tisi2009hot},
models@run.time~\cite{blair2009runtime}, and termination results for
transformation networks~\cite{stevens2020networks,klare2023termination}. Work
on LLM-generated domain models and model completion studies the quality of
single generated
models~\cite{chen2023domainmodeling,chaaben2023fewshot,clariso2023mdpe}; we
check how generated models and transformations fit together, a
well-formedness problem for transformation networks whose authors are LLMs and
whose values are stochastic.

% ===========================================================================
\section{Evaluation}\label{sec:eval}
% ===========================================================================

We ask three research questions, one per claim of the problem statement.
\textbf{RQ1 (Diagnosis):} do automatically assembled teams lack interfaces,
and how often is the decisive cause of their failures a composition defect?
\textbf{RQ2 (Prevention):} does the checker detect composition defects, and
how do typed, checked teams compare with free-form teams in task success,
reliability of completion, and cost?
\textbf{RQ3 (Attribution and repair):} when a typed team fails, does
trace-based attribution identify the responsible binding, agent and fault
class more accurately than transcript-based attribution, and does a checked
in-place repair preserve work?

\subsection{Setup}\label{sec:setup}

\paragraph*{Benchmarks} \emph{ClassEval}~\cite{du2024classeval} (100
class-level Python tasks, 4.1 methods per class on average, hidden unit tests
per method) is our primary benchmark: a class decomposes naturally into
per-method hand-offs, so composition matters. A task succeeds if all hidden
tests pass; we also report the per-test pass rate as a sensitive secondary
metric. \emph{HumanEval+}~\cite{chen2021codex,liu2023evalplus} (164
single-function tasks with extended hidden tests) is a control on which a
team should not help~\cite{kim2025scaling}. Teams, validators and builders
see only the task's skeleton and docstrings; hidden tests are used only for
scoring. Run on our sandbox, the reference solutions pass
\nCanonCE/100 ClassEval and \nCanonHE/164 HumanEval+ tasks (a few ClassEval
tests are time- or randomness-dependent).

\paragraph*{Models} We use two open-weight models served locally by Ollama,
for every role in every condition (builder, agents, critic, attribution):
\emph{Qwen2.5-Coder-7B}~\cite{hui2024qwen25coder} and
\emph{Qwen3.8-27B}~\cite{qwen2025qwen3} (thinking disabled), at temperature
0.6, a 16k context, and a cap of \nBudget{} generated tokens per run, equal
for all conditions.

\paragraph*{Conditions} We re-implemented the alternatives in one harness, so
that model, tools, worked example and budget are identical; each gets the
same execution tool, which runs the public docstring examples and any tests
the agents write. \textbf{Single}: one agent writes the solution and gets two
rounds of execution feedback. \textbf{Free}: a CaptainAgent-style
builder~\cite{song2024captain} writes 2--4 prose roles and a plan; the agents
talk in a round-robin group chat of at most 8 turns, which ends when an agent
says \code{TERMINATE}. \textbf{Critic}: Free, plus an LLM critic that reviews
the team for \D1--\D5 before it runs, and one revision. \textbf{Schema}: the
same roles exchange JSON patches to one shared board validated against a
schema, in the style of PatchBoard~\cite{zhang2026patchboard}.
\textbf{Typed-unchecked}: the \auto{} builder's first typed team, run
without the checker and without repair, which separates the format from the
check. \textbf{\auto{}}: Algorithms~\ref{alg:admit} and~\ref{alg:run} with
$k_{\mathit{adm}}=3$, $k=3$, $k_{\mathit{rep}}=2$, $r=1$.
\textbf{Typed-ref}: the hand-written admitted reference team (Tester and
Developer per method) run by Algorithm~\ref{alg:run}, which separates the
builder from the runtime. Every builder gets one worked example for the same
out-of-benchmark task (a \code{ParkingMeter} class), in its own format.

\paragraph*{Runs and statistics} \nRunsTotal{} runs in total. The 7B model runs
every condition on all ClassEval tasks with three seeds and on all
HumanEval+ tasks with two; the 27B model, about five times slower on our
hardware, runs one seed on a fixed subset of \nCESubset{} ClassEval and
\nHESubset{} HumanEval+ tasks (every fifth/fourth task, chosen before any
result was seen). We average success over seeds per task and compare
conditions per task with Wilcoxon signed-rank tests (Holm-corrected within
each benchmark--model cell), report exact McNemar counts on seed~1, paired
bootstrap 95\% confidence intervals, and Cliff's
$\delta$~\cite{arcuri2014guide}.

\todo{results}

% ===========================================================================
\section{Conclusion}\label{sec:conclusion}
% ===========================================================================
\todo{conclusion}

\noindent\textbf{Data availability.}\label{sec:data} \todo{artifact}

\bibliographystyle{IEEEtran}
\bibliography{references}

\end{document}
